\chapter{Planar Graphs}
%%%%%%%%%%%%%%%%%%%%%%%%%%%%%%%%%%%%%%%%%%%%%%%%%%%%%%%%%%%%%%%%%%%%%%%%%%%%%%%
\section{Planar Graphs}

% Generated by scripts/blueprint_restate.py from TCSlib.GraphTheory.Core.PlanarGraphs.
% Each body is the STATEMENT only, taken from the declaration's Lean docstring:
% the one-line summary plus the verbatim book statement.  Proof, proof plan,
% significance commentary and formalisation notes are excluded by construction.

\begin{definition}[A walk is internally disjoint from a subgraph $H$: no \emph{internal} vertex lies in $H$]
\lean{SimpleGraph.Walk.InternallyDisjointFrom}
\label{SimpleGraph.Walk.InternallyDisjointFrom}
\leanok
A walk is \textbf{internally disjoint from} a subgraph \texttt{H}: no \emph{internal}
vertex lies in \texttt{H}.

\emph{\texttt{W} is internally-disjoint from \texttt{H}
(that is, no internal vertex of \texttt{W} is a vertex of \texttt{H}).}
\end{definition}

\begin{definition}[B\&M's relation $~$ on $E(G) \ E(H)$: $e_{1} ~ e_{2}$ iff joined by a walk internally disjoint from $H$\dots]
\lean{SimpleGraph.bridgeRel}
\label{SimpleGraph.bridgeRel}
\leanok
\uses{SimpleGraph.Walk.InternallyDisjointFrom}
B\&M's relation \texttt{\textasciitilde{}} on \texttt{E(G) \textbackslash{} E(H)}:
\texttt{e$_{1}$ \textasciitilde{} e$_{2}$} iff joined by a walk internally disjoint from
\texttt{H} whose first and last edges are \texttt{e$_{1}$}, \texttt{e$_{2}$}.
\texttt{Relation.EqvGen} takes the reflexive-transitive-symmetric closure (discharging
"\texttt{\textasciitilde{}} is an equivalence relation").

\emph{Let \texttt{H} be a given subgraph of a graph
\texttt{G}.  We define a relation \texttt{$\sim$} on \texttt{E(G) \textbackslash{} E(H)} by the condition that \texttt{e$_{1}$ $\sim$ e$_{2}$} if
there exists a walk \texttt{W} such that (i) the first and last edges of \texttt{W} are
\texttt{e$_{1}$} and
\texttt{e$_{2}$}, respectively, and (ii) \texttt{W} is internally-disjoint from \texttt{H}.  It is easy to
verify that \texttt{$\sim$} is an equivalence relation on \texttt{E(G) \textbackslash{}
E(H)}.}
\end{definition}

\begin{definition}[The bridge of $H$ in $G$ containing the edge $e$: the subgraph induced by $e$'s $~$-class]
\lean{SimpleGraph.bridgeOf}
\label{SimpleGraph.bridgeOf}
\leanok
\uses{SimpleGraph.bridgeRel}
The \textbf{bridge} of \texttt{H} in \texttt{G} containing the edge \texttt{e}: the
subgraph induced by \texttt{e}'s \texttt{\textasciitilde{}}-class.

\emph{A subgraph of \texttt{G - E(H)} induced by an
equivalence class under the relation \texttt{$\sim$} is called a bridge of \texttt{H} in
\texttt{G}. It
follows immediately from the definition that if \texttt{B} is a bridge of \texttt{H},
then \texttt{B} is
a connected graph and, moreover, that any two vertices of \texttt{B} are connected by a
path that is internally-disjoint from \texttt{H}. It is also easy to see that two
bridges
of \texttt{H} have no vertices in common except, possibly, for vertices of \texttt{H}.}

The section then narrows: \emph{In this section we are concerned with the study of
bridges of a cycle \texttt{C}. Thus, to avoid repetition, we shall abbreviate "bridge of
\texttt{C}" to "bridge" in the coming discussion.}
\end{definition}

\begin{definition}[Vertices of attachment of the bridge of $e$ to the cycle traced by $c$: $V(B) \cap V(C)$]
\lean{SimpleGraph.attach}
\label{SimpleGraph.attach}
\leanok
\uses{SimpleGraph.bridgeOf}
Vertices of \textbf{attachment} of the bridge of \texttt{e} to the cycle traced by
\texttt{c}: \texttt{V(B) $\cap$ V(C)}. (B\&M's \texttt{V(B,H)}.)

\emph{For a bridge \texttt{B} of \texttt{H}, we write
\texttt{V(B) $\cap$ V(H) = V(B, H)}, and call the vertices in this set the vertices of
attachment of \texttt{B} to \texttt{H}.}

\emph{In a connected graph every bridge has at least one vertex of attachment, and in a
block every bridge has at least two vertices of attachment.  A bridge with \texttt{k}
vertices of attachment is called a \texttt{k}-bridge. Two \texttt{k}-bridges with the
same
vertices of attachment are equivalent \texttt{k}-bridges.}
\end{definition}

\begin{definition}[Position of a vertex along a cycle (index into its support)]
\lean{SimpleGraph.Walk.cycleIdx}
\label{SimpleGraph.Walk.cycleIdx}
\leanok
Position of a vertex along a cycle (index into its support).
\end{definition}

\begin{definition}[Four vertices in cyclic order on $c$]
\lean{SimpleGraph.Walk.InCyclicOrder}
\label{SimpleGraph.Walk.InCyclicOrder}
\leanok
\uses{SimpleGraph.Walk.cycleIdx}
Four vertices in cyclic order on \texttt{c}.
\end{definition}

\begin{definition}[Forward distance along $c$ from $a$ to $x$, in steps, wrapping at the base point]
\lean{SimpleGraph.Walk.cycleDist}
\label{SimpleGraph.Walk.cycleDist}
\leanok
\uses{SimpleGraph.Walk.cycleIdx}
Forward distance along \texttt{c} from \texttt{a} to \texttt{x}, in steps, wrapping at
the base point.
\end{definition}

\begin{definition}[$x$ lies on the closed arc of $c$ running forwards from $a$ to $b$]
\lean{SimpleGraph.Walk.OnArc}
\label{SimpleGraph.Walk.OnArc}
\leanok
\uses{SimpleGraph.Walk.cycleDist}
\texttt{x} lies on the closed arc of \texttt{c} running \textbf{forwards} from \texttt{a} to \texttt{b}.
\end{definition}

\begin{definition}[$a$ and $b$ are consecutive vertices of attachment of the bridge of $e$: both are attachments, and no\dots]
\lean{SimpleGraph.ConsecutiveAttach}
\label{SimpleGraph.ConsecutiveAttach}
\leanok
\uses{SimpleGraph.Walk.OnArc, SimpleGraph.attach}
\texttt{a} and \texttt{b} are \textbf{consecutive} vertices of attachment of the bridge of \texttt{e}: both
are attachments, and no other attachment lies on the arc between them.
\end{definition}

\begin{definition}[Two bridges of a cycle avoid one another: all the attachments of one lie in a single segment of the other]
\lean{SimpleGraph.Avoids}
\label{SimpleGraph.Avoids}
\leanok
\uses{SimpleGraph.ConsecutiveAttach, SimpleGraph.Walk.OnArc, SimpleGraph.attach}
Two bridges of a cycle \textbf{avoid} one another: all the attachments of one lie in a
single segment of the other.

\emph{Two bridges avoid one another if all the
vertices of attachment of one bridge lie in a single segment of the other bridge;
otherwise they overlap.}
\end{definition}

\begin{definition}[Two bridges of a cycle overlap: they do not avoid one another]
\lean{SimpleGraph.Overlaps}
\label{SimpleGraph.Overlaps}
\uses{SimpleGraph.Avoids}
Two bridges of a cycle \textbf{overlap}: they do not avoid one another.

\emph{Two bridges avoid one another if all the
vertices of attachment of one bridge lie in a single segment of the other bridge;
otherwise they overlap.}
\end{definition}

\begin{definition}[Two bridges of a cycle are skew: they have attachments $a, a', b, b'$ occurring in that cyclic order on\dots]
\lean{SimpleGraph.Skew}
\label{SimpleGraph.Skew}
\leanok
\uses{SimpleGraph.Walk.InCyclicOrder, SimpleGraph.attach}
Two bridges of a cycle are \textbf{skew}: they have attachments \texttt{a, a', b, b'}
occurring in
that cyclic order on \texttt{c} (\texttt{a, b} of one bridge, \texttt{a', b'} of the
other).

\emph{Two bridges \texttt{B} and \texttt{B'} are skew if
there are four distinct vertices \texttt{u}, \texttt{v}, \texttt{u'} and \texttt{v'} of
\texttt{C} such that \texttt{u} and
\texttt{v} are vertices of attachment of \texttt{B}, \texttt{u'} and \texttt{v'} are vertices of attachment of
\texttt{B'}, and the four vertices appear in the cyclic order \texttt{u}, \texttt{u'}, \texttt{v}, \texttt{v'} on
\texttt{C}.}
\end{definition}

\begin{definition}[A longest cycle]
\lean{SimpleGraph.Walk.IsLongestCycle}
\label{SimpleGraph.Walk.IsLongestCycle}
\leanok
A \textbf{longest cycle}.
\end{definition}

\begin{theorem}[Minimum degree at most five under the planar edge bound]
\lean{SimpleGraph.minDegree_le_five_of_card_edge_le}
\label{SimpleGraph.minDegree_le_five_of_card_edge_le}
\leanok
Let $G$ be a simple graph on a finite vertex set with $n \ge 3$ vertices, and write $E$
for its number of edges. If $E + 6 \le 3n$ — equivalently $E \le 3n - 6$ — then $G$ has
a vertex of degree at most $5$; that is, the minimum degree of $G$ satisfies $\delta(G)
\le 5$.
\end{theorem}

\begin{theorem}[Edge counts of a graph and its complement]
\lean{SimpleGraph.card_edgeFinset_add_card_edgeFinset_compl}
\label{SimpleGraph.card_edgeFinset_add_card_edgeFinset_compl}
\leanok
Let $G$ be a simple graph on a finite vertex set $V$ with $n = \abs{V}$ vertices, and
let $G^{c}$ denote its complement, the simple graph on $V$ in which two distinct
vertices are adjacent precisely when they are not adjacent in $G$. Then the number of
edges of $G$ together with the number of edges of $G^{c}$ equals the number of unordered
pairs of vertices, \[ \abs{E(G)} + \abs{E(G^{c})} = \binom{n}{2}. \]
\end{theorem}

\begin{theorem}[Complement cannot also meet the edge bound at eleven vertices]
\lean{SimpleGraph.not_both_card_edge_le_of_eleven_le}
\label{SimpleGraph.not_both_card_edge_le_of_eleven_le}
\leanok
Let $G$ be a simple graph on a finite vertex set $V$ with $n = \abs{V} \ge 11$, write
$E(G)$ for its edge set, and let $G^{c}$ be its complement (the simple graph on $V$ in
which two distinct vertices are adjacent exactly when they are not adjacent in $G$).
Then $G$ and $G^{c}$ cannot both satisfy $\abs{E(G)} + 6 \le 3n$ and $\abs{E(G^{c})} + 6
\le 3n$; that is, at least one of the two graphs has more than $3n - 6$ edges.
\end{theorem}

\begin{theorem}[Distinct bridges meet only in the subgraph]
\lean{SimpleGraph.bridge_inter_subset_cycle}
\label{SimpleGraph.bridge_inter_subset_cycle}
\leanok
\uses{SimpleGraph.bridgeOf, SimpleGraph.bridgeRel}
Let $H$ be a subgraph of a finite simple graph $G$, and let $e$ and $e'$ be edges of $G$
that do not belong to $H$, so both lie in $E(G)\setminus E(H)$. Recall that two such
edges are called equivalent when they are joined by a walk, internally disjoint from
$H$, whose first and last edges are the two given ones, and that the bridge of $H$
containing an edge $f$ is the subgraph of $G$ induced by the equivalence class of $f$.
Write $B$ for the bridge containing $e$ and $B'$ for the bridge containing $e'$. If $e$
and $e'$ are not equivalent, then every vertex lying in both $B$ and $B'$ is a vertex of
$H$; that is, $V(B)\cap V(B')\subseteq V(H)$.
\end{theorem}

\begin{theorem}[Connectedness of a bridge of a subgraph]
\lean{SimpleGraph.bridgeOf_connected}
\label{SimpleGraph.bridgeOf_connected}
\leanok
\uses{SimpleGraph.bridgeOf}
Let $G$ be a finite simple graph, let $H$ be a subgraph of $G$, and let $e$ be an edge
of $G$ that does not belong to $H$. Then the bridge of $H$ in $G$ containing $e$ — the
subgraph induced by the $\sim$-equivalence class of $e$ — is connected: it has at least
one vertex, and any two of its vertices are joined by a path within it.
\end{theorem}

\begin{theorem}[Vertices of a bridge are joined by internally disjoint paths]
\lean{SimpleGraph.exists_path_internallyDisjoint}
\label{SimpleGraph.exists_path_internallyDisjoint}
\leanok
\uses{SimpleGraph.Walk.InternallyDisjointFrom, SimpleGraph.bridgeOf}
Let $G$ be a finite simple graph, let $H$ be a subgraph of $G$, and let $e$ be an edge
of $G$ that does not belong to $H$. Then any two vertices of the bridge of $H$ in $G$
containing $e$ are joined by a path in $G$ that is internally disjoint from $H$.
\end{theorem}

\begin{theorem}[Tripod at a bridge with three vertices of attachment]
\lean{SimpleGraph.bridge_tripod_of_three_attachments}
\label{SimpleGraph.bridge_tripod_of_three_attachments}
\leanok
\uses{SimpleGraph.bridgeOf}
Let $C$ be a cycle in a finite simple graph $G$, and let $e$ be an edge of $G$ that does
not lie on $C$; write $B$ for the bridge of $C$ in $G$ that contains $e$. Suppose $v_1,
v_2, v_3$ are three distinct vertices of attachment of $B$, that is, three distinct
vertices each lying both in $B$ and on $C$. Then there exists a vertex $v_0$ of $B$ that
does not lie on $C$, together with three paths $P_1, P_2, P_3$ in $G$ joining $v_0$ to
$v_1$, $v_2$, and $v_3$ respectively, such that any two of these paths have only the
vertex $v_0$ in common.
\end{theorem}

\begin{theorem}[Overlapping bridges are skew or equivalent 3-bridges]
\lean{SimpleGraph.overlap_imp_skew_or_equivalent_three_bridge}
\label{SimpleGraph.overlap_imp_skew_or_equivalent_three_bridge}
\leanok
\uses{SimpleGraph.Overlaps, SimpleGraph.Skew, SimpleGraph.attach}
Let $G$ be a finite simple graph, let $C$ be a cycle in $G$, and let $e$ and $e'$ be
edges of $G$ that do not lie on $C$; write $B$ and $B'$ for the bridges of $C$
containing $e$ and $e'$ respectively. If $B$ and $B'$ overlap, then either they are
skew, or their sets of vertices of attachment to $C$ coincide and consist of exactly
three vertices — that is, $B$ and $B'$ are equivalent $3$-bridges.
\end{theorem}

\begin{theorem}[A bridge of a longest cycle reaching off the cycle]
\lean{SimpleGraph.exists_bridge_off_longest_cycle}
\label{SimpleGraph.exists_bridge_off_longest_cycle}
\leanok
\uses{SimpleGraph.Walk.IsLongestCycle, SimpleGraph.bridgeOf}
Let $G$ be a finite connected simple graph that is not Hamiltonian, and let $C$ be a
longest cycle in $G$ — a cycle whose length is at least that of every cycle of $G$. Then
there is an edge $e \in E(G) \setminus E(C)$ such that the bridge of $C$ containing $e$
has a vertex that does not lie on $C$.
\end{theorem}

\begin{theorem}[Non-adjacency of successors of a bridge's attachment vertices on a longest cycle]
\lean{SimpleGraph.succ_attachments_not_adj}
\label{SimpleGraph.succ_attachments_not_adj}
\leanok
\uses{SimpleGraph.Walk.IsLongestCycle, SimpleGraph.attach, SimpleGraph.bridgeOf}
Let $G$ be a finite connected graph, and let $C$ be a longest cycle of $G$, with its
vertices listed in cyclic order as $v_0, v_1, v_2, \dots$ (so $v_k$ is the vertex
reached after $k$ steps along $C$, and $v_{k+1}$ is its successor). Let $e$ be an edge
of $G$ that does not lie on $C$, and let $B$ be the bridge of $C$ that contains $e$;
suppose $B$ has at least one vertex not lying on $C$. If $v_i$ and $v_j$ are two
distinct vertices of attachment of $B$ to $C$, then $v_{i+1}$ and $v_{j+1}$ are not
adjacent in $G$.
\end{theorem}

\begin{theorem}[The Chvátal–Erdős condition for Hamiltonicity]
\lean{SimpleGraph.chvatal_erdos_hamiltonian}
\label{SimpleGraph.chvatal_erdos_hamiltonian}
\leanok
\uses{SimpleGraph.vertexConnectivity}
Let $G$ be a finite simple graph on at least three vertices, and suppose $G$ is
connected. Write $\alpha(G)$ for the independence number of $G$ (the largest size of a
set of pairwise nonadjacent vertices) and $\kappa(G)$ for its vertex connectivity. If
$\alpha(G) \le \kappa(G)$, then $G$ is Hamiltonian; that is, $G$ has a cycle passing
through every vertex exactly once.
\end{theorem}

\begin{theorem}[Tait colouring of Hamiltonian cubic graphs]
\lean{SimpleGraph.lineGraph_colorable_three_of_isHamiltonian_of_three_regular}
\label{SimpleGraph.lineGraph_colorable_three_of_isHamiltonian_of_three_regular}
\leanok
Let $G$ be a finite simple graph that is $3$-regular, so every vertex has exactly three
neighbours, and suppose $G$ admits a Hamiltonian cycle, a closed walk that visits every
vertex exactly once. Then the line graph of $G$ — whose vertices are the edges of $G$,
two being adjacent precisely when they share an endpoint — is $3$-colourable;
equivalently, the edges of $G$ can be properly coloured with three colours, a Tait
colouring.
\end{theorem}
